% =====================================================================
% Section 1: Introduction — EvidenceFlow
% =====================================================================
%
% Status: Draft §1 Introduction v4 (2026-05-17)
% Revision note (v4 vs v3):
%   - Contribution bullets finalized after §4 (fact-source substrate +
%     local PPR + Evidence-Aware Retrieval Optimization) frozen.
%   - Bullet 2 now exposes the fact-source substrate explicitly
%     (matching §4.1's central claim) so the Method's distinguishing
%     feature is visible in the contribution list.
%   - Bullet 3 names the three datasets (HotpotQA / 2Wiki / MuSiQue)
%     to match the evaluation paragraph and avoid generic "three
%     benchmarks".
%   - Removed the headline-numbers and contributions-finalization
%     TODO markers; numbers go into abstract + final bullet at
%     submission once the main table is frozen.
%
% Revision note (v3 vs v2): see archive/ — de-duplicated §1 vs §2,
% replaced generic opening, anchored with named failure modes.
%
% =====================================================================

\section{Introduction}
\label{sec:intro}

Multi-hop question answering forces a retriever to deliver evidence
whose value depends on relations across documents, not on individual
question-passage similarity. A useful retrieval result must contain both
the bridge passages that link question entities to the answer and the
tail passages that resolve each entity or relation the question
mentions. Pointwise dense retrieval was designed for the opposite
regime, in which one passage answers one question, and tends to
concentrate the top-$K$ list around the question's most lexically
salient entity while leaving multi-hop bridges scored too low to
participate.

Graph-based retrieval-augmented generation, or GraphRAG, addresses this
mismatch by organizing the corpus into a structure that exposes
cross-document connections and biasing retrieval toward that structure
rather than toward surface similarity. Existing GraphRAG systems differ
mainly in how they introduce structure. One line builds an
OpenIE-derived graph and propagates relevance from question seeds
through personalized PageRank, with
HippoRAG~\citep{jimenez2024hipporag,gutierrez2025hipporag} as a
representative instance. A second line makes multi-hop chains explicit
by searching or constructing reasoning paths at query time, exemplified
by PropRAG~\citep{wang2025proprag}. A third line replaces the underlying
representation, using community summaries or hypergraphs over the
corpus, as in GraphRAG~\citep{edge2024graphrag} and
HGRAG~\citep{hgrag2024}. We discuss these systems and their relation to
\methodname{} in \S\ref{sec:related}.

These routes all face the same downstream constraint. Whatever the
graph operator returns, retrieval-first multi-hop QA still has to hand
the reader a short ranked list of natural-language passages under a
small budget. Two failure modes recur in this conversion. Propagation
puts mass on the strongest seed neighborhood, so a top-$K$ list can
look on-topic for the question and still omit the bridge passages that
connect to the answer. Independently, a list of high-scoring passages
can crowd around a single entity in the question while leaving other
entities or relations entirely uncovered, so the reader has the
ingredients for one hop and nothing for the next. Path search and
summary representations can shift where this conversion is handled, but
they still leave the retriever with a final reader-facing ordering
decision.

The problem is therefore not only how to propagate relevance over a
graph, but how to turn propagated relevance into evidence that is both
\emph{connected}---preserving the bridge relations that span
documents---and \emph{covering}, in the sense that the selected
passages collectively address the distinct question-side entities and
relations a reader has to resolve. We treat this conversion as the core
retrieval problem and keep the interface simple. The system outputs a
ranked evidence list $R_q$ of natural-language passages, not a symbolic
path or a generated reasoning chain, and the reader consumes a prefix
of $R_q$.

We propose \methodname{}, a graph-based retrieval framework built on a
single design principle. Graph propagation is the start of multi-hop
evidence selection, not the final ranking. \methodname{} operationalizes
this principle through a staged retrieval pipeline. Offline indexing
builds a document-grounded OpenIE graph that
couples passages, entity phrases, facts, and provenance
(\S\ref{sec:method:substrate}). At query time, fact-level personalized
PageRank on a query-local view of this graph yields an evidence flow
$\pi_q$ (\S\ref{sec:method:retrieval}). Because $\pi_q$ can underweight
bridges several hops from any seed, \methodname{} then follows
relation-grounded edges to admit bridge and tail-entity passages, and
reranks the merged candidate set with a noisy-OR coverage model over
question-side requirements (\S\ref{sec:method:enhancement}). The
query-local propagation stage keeps retrieval focused,
relation-grounded expansion repairs missing bridges, and coverage-aware
reranking prevents the final list from collapsing into a single entity
neighborhood.

We evaluate \methodname{} on HotpotQA, 2WikiMultiHopQA, and MuSiQue
under a common reader, comparing against dense, graph-based, and
path-oriented retrieval baselines. The evaluation reports retrieval
quality and downstream QA performance, and ablates each of the three
stages. The experiments are designed to test whether the proposed
conversion---from graph relevance to a connected and covering evidence
list---improves both support coverage and answer quality.

The main contributions are as follows.
\begin{itemize}
    \item We frame retrieval-first multi-hop QA as a conversion
    problem of turning graph-level relevance into a connected and
    covering evidence list under a small reader budget, and we
    identify two failure modes (bridge underweighting and
    entity-neighborhood crowding) this conversion has to address.
    \item We instantiate this framing in \methodname{} on top of a
    fact-source substrate that promotes each OpenIE fact to a graph
    node and binds it to its source passage through endpoint and
    source edges. On this substrate, query-local fact-level
    personalized PageRank, relation-grounded expansion, and
    noisy-OR coverage reranking jointly target focused propagation,
    bridge recovery, and requirement coverage.
    \item We show on HotpotQA, 2WikiMultiHopQA, and MuSiQue that
    this design consistently improves retrieval and answer quality
    over dense, graph-based, and path-oriented retrieval baselines
    under a unified reader, with component ablations isolating the
    contribution of each operator.
\end{itemize}
